\chapter{Problem Formulation}

\section{Formal Problem Specification}

\section{Experimental Set-ups}

\subsection{Summary}

To begin the dissertation, I examinied the simplest forms of automatic fact verification for table-based fact verification. These are methods whose input is the relevant table as gold evidence concatenated with the claim to verify, but no relevant rows identified. From that input the models have to determine if the claim is entailed or refuted. The models can be used as baselines to assess single table classification performance, which can be compared to their performance on a multi-table fact verificaion dataset design and be useful to inform a different pipeline for automatic verification of table-based data I could develop later. I worked with the TabFact dataset \parencite{chen2019tabfact}, which is a widely used standard for table-based fact verification with a balanced dataset of entailed and refuted claims, as well as diversity in the complexity of logic to reach a conclusion. I conducted a thorough analysis of the dataset beforehand. All the work to read in the TabFact datasets, analyse data, do supervised fine-tuning, and perform classification was done in Python with a Google Colab Jupyter notebook utilising cloud CPUs and GPUs. For the language model baselines, I used Hugging Face transformers \parencite{wolf2020transformers} and PyTorch \parencite{paszke2019pytorch}.

\subsection{Inputs for Single Table Baselines} 
The input for each singel table pipeline is the table and the claim concatenated together. This was done with the TabFact linearised table straight from the dataset coming after “TABLE: ”  and the claim to be verified coming after “CLAIM: ”. For each baseline, I made a reasonable effort to produce a result that was not a single-class prediction to estimate a reasonable ceiling for performance. Therefore I varied the hyperparameters slightly including the input of the table and the claim itself and observed validation set performance. Sometimes I would vary the separation of the table claim with concatenation with white space and a pipe character, or I would simply start a newline. Specific adjustments and training configurations for each baseline can be found in the appendix. A typical table before claim (T+C) and claim before table (C+T) input is displayed in Figure~\ref{fig:input-formats}.
\begin{figure}[htbp]
\centering
\small
\begin{tabular}{|p{0.94\linewidth}|}
\hline
\textbf{(a) Table-before-Claim (T+C)} \\
\ttfamily\footnotesize
TABLE: date\#label\#region\#format\#catalog \\
march 1984\#rca\#france\#7 single (first cover)\#61298 \\
september 1984\#rca\#france\#7 single (second cover)\#61298 \\
september 1984\#rca\#france\#7 maxi\#61299 \\
september 1984\#rca\#france\#7 single (english version)\#61531 \\
1985\#original\#canada\#7 single\#6502 \\
CLAIM: the second cover of catalog 61298 was released by rca on september 1984 \\
\hline
\textbf{(b) Claim-before-Table (C+T)} \\
\ttfamily\footnotesize
CLAIM: the athlete or team with the rank of 3 won 3 gold , 2 silver , and 3 bronze medals \\
TABLE: rank\#gold\#silver\#bronze\#total \\
1\#7\#4\#2\#13 \\
2\#3\#4\#4\#11 \\
3\#3\#2\#3\#8 \\
4\#1\#0\#0\#1 \\
5\#0\#2\#0\#2 \\
6\#0\#1\#3\#4 \\
7\#0\#1\#2\#3 \\
total\#14\#14\#14\#42 \\
\hline
\end{tabular}
\caption{Examples of the linearised input formats used for baseline experiments, illustrating (a) Table-before-Claim (T+C) and (b) Claim-before-Table (C+T) orderings.}
\label{fig:input-formats}
\end{figure}

\subsection{Baseline Selection}
The baselines were chosen based on precedence in text and fact verification tasks as a method to understand the problem of table and claim fact verification and provide simple pipelines which I can compare when I evaluate performance in multi-table fact verification pipelines and a different approach I design myself later in the dissertation. Logistic regression \parencite{cox1958regression} is a good standard practical choice in NLI as it is quick. By attaining simple vector representations of the table and claim together, we can establish if pure linguistic features of the table and claim together can map to an accurate classification. We do not expect it to perform well, because logistic regression is not pretrained and is unable to perform semantic understanding of the table structure, showing the need to use language models. The language models BERT \parencite{devlin2019bert}, GPT2 \parencite{radford2019language}, and Qwen3-0.6B \parencite{yang2025qwen3}, were chosen, as they all have strong precedent for use in text classification problems which require semantic understanding, including structured data and for fact verification, with specific examples being cited in the literature survey. 

\subsubsection{Logistic Regression}
\noindent For logistic regression, the table claim concatenation was converted into a TFIDF vector. I used scikit-learn’s \parencite{pedregosa2011scikit} LogisticRegression implementation to fit the training data and then generate predictions on the test set.

\subsubsection{Language Models}
\noindent For the sequence classifications using BERT, GPT-2, and Qwen3-0.6B, I used their respective tokenisers to preprocess the input data and then used the corresponding HuggingFace \{model name\}ForSequenceClassification class with the Trainer class to train the model and generate predictions. There were truncation issues with most of the tokenisers, as some of the tables were very long, and for practicality, some had to be truncated more than others. Detailed accounts of this, as well as training parameters, are mentioned in the appendix. As Qwen3-0.6B is too large a model for me to practically fine-tune, I used Low-Rank Adaptation (Hu et al., 2021) for parameter-efficient fine-tuning (PEFT), which freezes Qwen’s base model weights and updates compact low-rank adapter parameters instead.

\subsubsection{Additional Qwen3 focus}
\noindent For Qwen3-0.6B, I also tested standard zero-shot and few-shot testing performance \parencite{brown2020language}. These were all conducted with thinking mode disabled, but I also included a randomised sample of the testing dataset with thinking mode enabled to gain an additional observation on the capabilities of Qwen's explicit reasoning for fact verification. The pursuit of so many baselines for Qwen compared to other models comes down to the fact it is a new model with a lot of recent interest in its text classification potential but not many papers out currently.

\subsection{Evaluation Metrics}
\noindent The evaluation metrics chosen were accuracy, precision, recall, and F1 score \parencite{sokolova2009systematic} with respect to the positive entailed class. These were used in the TabFact paper, with accuracy as the primary evaluation metric, so that and the dataset being balanced is why it is my primary metric also. These are also just standard models of evaluation with natural language inference generally \parencite{bowman2015large} and other fact-checking benchmarks \parencite{thorne2018fever}. The secondary metrics are valuable for spotting single-class prediction biases.

\subsection{Limitations}
Time and resource constraints and advice from my supervisor justified using a headline test result from a single seed to represent baseline success for the exploratory baselines. Where possible, model validation scores such as in logistic regression, BERT, and GPT-2 were run, and means and standard deviations were calculated to confirm that the models were stable on their validation sets. It was not feasible to do so for many of the Qwen classification results as they took an extremely long time to fine-tune, even though it would have been ideal to calculate them. That is a limitation I am willing to accept and doing a test for two variations of representing the table and claim are enough to show promise at this stage. 
